% !PRINTED 502
\weberpage{502}

Nach § 38 (15) ist also

\weq{10}{M = (-1)^{A + C + \frac{B+1}{2}} \sqrt{D}^3\, \gamma_3(\omega),}
und es wird hier nur einer der Werte $v = u$, demnach ist $M$
und folglich auch $\sqrt{D}\,\gamma_3(\omega)$ rational durch $j(\omega)$ ausdrückbar.

Um auch den Ausnahmefall zu erledigen, setzen wir
\[
n = -D = 3\,m^2.
\]
Dann wird (9):
\[
y^2 + 3\,m^2(x^2 - 4) = 0,
\]
und diese hat die drei Lösungen:
\[
\begin{aligned}
x &= 2, & y &= 0, \\
x &= 1, & y &= \mp 3\,m.
\end{aligned}
\]
Die Gleichungen (8) ergeben für $x = 2$ wieder die Formel (10):
\[
M_1 = (-1)^{A + C + \frac{B+1}{2}} \sqrt{D}^3\, \gamma_3(\omega),
\]
und die beiden anderen Fälle ergeben
\[
\begin{aligned}
\partial &= 1, & a &= n, \\
\alpha &= \frac{B \pm 3\,m}{2}, & \beta &= A, \\
\gamma &\equiv C, & \alpha + \delta &\equiv B \pmod{4}, \\
\alpha + \beta\omega &= \frac{\sqrt{D} \pm 3\,m}{2}. &&
\end{aligned}
\]
Daraus erhält man nach § 38 (15) für $M_2$, $M_3$ die Werte
\[
M_2 = (-1)^{\frac{B-1}{2}} \left(\frac{3\,m + \sqrt{D}}{2}\right)^3 \gamma_3(\omega),
\]
\[
M_3 = (-1)^{\frac{B-1}{2}} \left(\frac{-3\,m + \sqrt{D}}{2}\right)^3 \gamma_3(\omega).
\]
Nun sind zwar nicht die Größen $M_1$, $M_2$, $M_3$ einzeln, wohl
aber, wie wir gleich zeigen werden, ihre symmetrischen Funktionen, z. B. ihre Summe, durch $j(\omega)$ rational ausdrückbar, und
daraus ergibt sich auch für diesen Fall, daß $\sqrt{D}\,\gamma_3(\omega)$ Klasseninvariante ist.

Von der Annahme, die wir gemacht haben, daß $A$ relativ
prim zu $2\,D$ sei, können wir uns nachträglich befreien, da $\gamma_3(\omega)$
durch alle linearen Transformatoren höchstens sein Zeichen
ändert. Wir haben dann den Satz:

1. \gesperrt{Ist $D \equiv 1 \pmod{4}$, so ist $\sqrt{D}\,\gamma_3(\omega)$ Klasseninvariante der Diskriminante $D$.}
